% ==============================================================================
% ANNEXE 2 -- GRILLE D'ANALYSE THÉMATIQUE (CODEBOOK)
% Mémoire MAE -- Arthur SAUNIER -- Wavestone Luxembourg
% Le \chapter{} est déclaré dans le fichier maître ; ce fichier démarre en \section.
% ==============================================================================

\section{Principes de construction de la grille}

L'analyse du matériau recueilli repose sur une analyse thématique de contenu conduite
selon une logique mixte. L'ossature de la grille est \textbf{déductive}~: les quatre
axes de codage sont directement issus des quatre questions directrices et de leurs
ancrages théoriques respectifs, ce qui autorise la confrontation systématique entre la
littérature et le terrain conduite au chapitre~4. Les codes de rang inférieur sont
quant à eux partiellement \textbf{inductifs}~: un axe transversal ouvert, désigné par
la lettre~T afin qu'aucun code ne se confonde avec un code de répondant, accueille les
thèmes non anticipés rencontrés à la lecture du corpus. La grille est appliquée
rétroactivement à l'ensemble des entretiens, selon un référentiel unique.

Les conventions de codage retenues sont les suivantes.

\begin{itemize}
    \item \textbf{Unité de codage.} L'unité retenue est le segment de sens, c'est-à-dire
    l'extrait de discours porteur d'une idée complète, indépendamment de sa longueur
    grammaticale. Un même segment peut recevoir plusieurs codes lorsqu'il articule
    plusieurs dimensions.

    \item \textbf{Codage de la contradiction.} Lorsqu'un répondant énonce une position
    contredite ailleurs dans le même entretien, les deux segments sont codés et
    l'écart est signalé par le code transversal T1.1. Ces contradictions ne sont pas
    arbitrées mais constituent un résultat en soi.

    \item \textbf{Traitement de la saturation.} La saturation est considérée comme
    atteinte lorsqu'un entretien supplémentaire n'ajoute aucun code nouveau et ne
    modifie pas la structure de la grille. Le rang de l'entretien à partir duquel
    aucun code inédit n'apparaît est mentionné dans la restitution.

    \item \textbf{Anonymisation.} Chaque répondant est désigné par un code de profil
    (J pour junior --- analyste ou consultant ---, E pour encadrant --- du consultant
    senior à l'associé), suivi
    d'un rang d'entretien. Les noms de clients et de missions sont remplacés par des
    désignations génériques entre crochets dans les verbatims cités.

    \item \textbf{Traçabilité.} Chaque verbatim reproduit dans le mémoire est référencé
    par le code du répondant et le code thématique, de manière à ce que tout lecteur
    puisse reconstituer le chemin analytique entre le corpus brut et l'interprétation
    proposée.
\end{itemize}

\section{Axe A --- L'appropriation des outils dans le travail d'analyse}

\textit{Adossé à la question directrice~1. Cadres mobilisés~: sociologie des usages,
modèle d'acceptation technologique (Davis, 1989).}

\begin{table}[h]
\centering
\small
\begin{tabular}{|p{0.08\linewidth}|p{0.27\linewidth}|p{0.51\linewidth}|}
\hline
\textbf{Code} & \textbf{Intitulé} & \textbf{Définition et critère d'affectation} \\
\hline
A1.1 & Recherche et veille & Délégation de la recherche documentaire, réglementaire ou sectorielle \\
\hline
A1.2 & Synthèse et reformulation & Condensation ou réécriture d'un matériau existant \\
\hline
A1.3 & Première rédaction & Production d'un premier jet de livrable ou de section \\
\hline
A1.4 & Traduction et mise en forme & Travail linguistique ou formel sans production analytique \\
\hline
A1.5 & Structuration du raisonnement & Aide à la construction d'un plan ou d'une problématique \\
\hline
A2.1 & Sanctuarisation --- confidentialité & Tâche non déléguée en raison de la sensibilité des données \\
\hline
A2.2 & Sanctuarisation --- jugement & Tâche non déléguée car jugée exiger un jugement professionnel \\
\hline
A2.3 & Sanctuarisation --- relation & Tâche non déléguée car relevant de l'interaction client \\
\hline
A3.1 & Apprentissage autonome & Montée en compétence sur l'outil par essai personnel \\
\hline
A3.2 & Transmission entre pairs & Apprentissage horizontal, informel, au sein de l'équipe \\
\hline
A3.3 & Formation institutionnelle & Dispositif de formation porté par le cabinet \\
\hline
A4.1 & Utilité perçue & Bénéfice attendu exprimé en termes de performance (TAM) \\
\hline
A4.2 & Facilité d'usage perçue & Effort d'usage exprimé comme faible ou élevé (TAM) \\
\hline
A4.3 & Résistance ou évitement & Refus, réticence ou abandon déclaré de l'outil \\
\hline
A4.4 & Contournement outillé & Recours à un outil personnel hors du cadre officiel \\
\hline
\end{tabular}
\caption{Axe A --- Codes relatifs à l'appropriation des outils}
\label{tab:codebook-a}
\end{table}

\section{Axe B --- Apprentissage, compétences et transmission}

\textit{Adossé à la question directrice~2. Cadres mobilisés~: Polanyi (1966) sur le
savoir tacite, Nonaka et Takeuchi (1995) sur la conversion des connaissances.}

\begin{table}[h]
\centering
\small
\begin{tabular}{|p{0.08\linewidth}|p{0.27\linewidth}|p{0.51\linewidth}|}
\hline
\textbf{Code} & \textbf{Intitulé} & \textbf{Définition et critère d'affectation} \\
\hline
B1.1 & Vertu formatrice reconnue & Tâche ingrate décrite comme pédagogiquement nécessaire \\
\hline
B1.2 & Vertu formatrice déniée & Tâche ingrate décrite comme sans apport d'apprentissage \\
\hline
B2.1 & Accélération de la montée en compétence & Progression jugée plus rapide grâce à l'outil \\
\hline
B2.2 & Atrophie du raisonnement & Perte déclarée ou observée de capacité d'analyse autonome \\
\hline
B2.3 & Vérification comme compétence & Émergence du contrôle critique comme savoir-faire propre \\
\hline
B3.1 & Relecture senior formatrice & La revue hiérarchique décrite comme moment de transmission \\
\hline
B3.2 & Appauvrissement de l'interaction & Réduction des échanges junior-senior imputée à l'outil \\
\hline
B3.3 & Compagnonnage préservé & Maintien affirmé des mécanismes traditionnels de socialisation \\
\hline
B4.1 & Dépendance déclarée & Difficulté exprimée à travailler sans l'outil \\
\hline
B4.2 & Validation sans maîtrise & Approbation d'un contenu que le répondant ne saurait produire \\
\hline
\end{tabular}
\caption{Axe B --- Codes relatifs à l'apprentissage et aux compétences}
\label{tab:codebook-b}
\end{table}

\section{Axe C --- Légitimité de l'expert et relation client}

\textit{Adossé à la question directrice~3. Cadres mobilisés~: Abbott (1988) sur la
juridiction professionnelle, Suchman (1995) sur les registres de légitimité.}

\begin{table}[h]
\centering
\small
\begin{tabular}{|p{0.08\linewidth}|p{0.27\linewidth}|p{0.51\linewidth}|}
\hline
\textbf{Code} & \textbf{Intitulé} & \textbf{Définition et critère d'affectation} \\
\hline
C1.1 & Légitimité cognitive & Confiance fondée sur l'expertise et la maîtrise du domaine \\
\hline
C1.2 & Légitimité pragmatique & Confiance fondée sur l'utilité et le résultat obtenu \\
\hline
C1.3 & Légitimité morale & Confiance fondée sur la responsabilité assumée et la déontologie \\
\hline
C2.1 & Transparence assumée & Usage de l'outil déclaré au client \\
\hline
C2.2 & Opacité stratégique & Usage délibérément non déclaré \\
\hline
C2.3 & Interpellation par le client & Le client soulève lui-même la question de l'usage \\
\hline
C3.1 & Déplacement vers le jugement & Valeur du cabinet redéfinie autour du sens critique \\
\hline
C3.2 & Déplacement vers la responsabilité & Valeur du cabinet redéfinie autour de l'engagement pris \\
\hline
C3.3 & Menace de désintermédiation & Crainte que le client produise l'analyse par lui-même \\
\hline
C4.1 & Tension sur le modèle tarifaire & Remise en cause de la facturation au temps passé \\
\hline
\end{tabular}
\caption{Axe C --- Codes relatifs à la légitimité et à la relation client}
\label{tab:codebook-c}
\end{table}

\section{Axe D --- Gouvernance, contrôle qualité et risques}

\textit{Adossé à la question directrice~4. Cadres mobilisés~: Mintzberg (1982) sur les
mécanismes de coordination de la bureaucratie professionnelle, gouvernance des risques
opérationnels.}

\begin{table}[h]
\centering
\small
\begin{tabular}{|p{0.08\linewidth}|p{0.27\linewidth}|p{0.51\linewidth}|}
\hline
\textbf{Code} & \textbf{Intitulé} & \textbf{Définition et critère d'affectation} \\
\hline
D1.1 & Règle formelle connue & Référence à une politique explicite du cabinet \\
\hline
D1.2 & Norme informelle & Règle d'équipe non écrite, transmise oralement \\
\hline
D1.3 & Flou normatif & Méconnaissance ou incertitude sur le cadre applicable \\
\hline
D2.1 & Revue hiérarchique & Contrôle du livrable par un niveau supérieur \\
\hline
D2.2 & Vérification des sources & Contrôle spécifique de la véracité du contenu produit \\
\hline
D2.3 & Absence de contrôle dédié & Aucun dispositif spécifique au contenu généré \\
\hline
D3.1 & Hallucination détectée & Incident avéré de contenu erroné identifié à temps \\
\hline
D3.2 & Quasi-incident & Erreur passée près d'atteindre le client \\
\hline
D3.3 & Enjeu de confidentialité & Risque de divulgation de données de mission \\
\hline
D4.1 & Attribution de responsabilité & Désignation de qui répond de l'erreur \\
\hline
D5.1 & Règle perçue comme bridante & Cadre jugé contraignant au regard du travail réel \\
\hline
D5.2 & Demande de cadrage & Souhait exprimé d'un encadrement plus explicite \\
\hline
\end{tabular}
\caption{Axe D --- Codes relatifs à la gouvernance et aux risques}
\label{tab:codebook-d}
\end{table}

\section{Axe T --- Codes transversaux et émergents}

\begin{table}[h]
\centering
\small
\begin{tabular}{|p{0.08\linewidth}|p{0.27\linewidth}|p{0.51\linewidth}|}
\hline
\textbf{Code} & \textbf{Intitulé} & \textbf{Définition et critère d'affectation} \\
\hline
T1.1 & Écart discours-pratique & Contradiction interne entre position déclarée et pratique décrite \\
\hline
T1.2 & Écart junior-encadrant & Divergence de perception entre les deux positions sur un même objet \\
\hline
T2.1 & Prospective du métier & Projection sur l'évolution du conseil à moyen terme \\
\hline
T2.2 & Recommandation spontanée & Proposition d'action formulée par le répondant \\
\hline
T3.1 & Le non-dit interne & Usage de l'outil passé sous silence entre collègues, indépendamment du client \\
\hline
T3.2 & L'incident déclencheur de contrôle & Modification d'une pratique de revue consécutive à une erreur vécue \\
\hline
\end{tabular}
\caption{Axe T --- Codes transversaux et catégorie ouverte}
\label{tab:codebook-t}
\end{table}

\section{Matrice de restitution}

Le corpus codé est reporté dans une matrice croisant les codes en lignes et les
répondants en colonnes, complétée d'une colonne de fréquence d'occurrence et d'une
colonne de verbatim représentatif. Cette matrice constitue le support direct de la
restitution du chapitre~4~: chaque sous-partie analytique s'y adosse en signalant les
codes saturés, les codes clivants entre profils et les codes isolés. La lecture en
colonnes restitue la cohérence interne du discours de chaque répondant~; la lecture en
lignes fait apparaître les convergences et les lignes de fracture au sein du panel,
notamment entre le groupe des juniors et celui des encadrants, dont la confrontation
constitue l'un des principaux résultats attendus de l'étude.